\chapter{Pipeline de traitement du signal et extraction des indicateurs}
\label{chap:pipeline}

\section{Prétraitement}
Tout signal entrant subit une chaîne de prétraitement homogène avant analyse :
\begin{enumerate}[leftmargin=1.6em]
  \item \textbf{Rééchantillonnage} à la fréquence cible $f_s = 12\,800$\,Hz par
        \texttt{resample\_poly} (SciPy), afin que tous les datasets partagent une même
        base temporelle.
  \item \textbf{Normalisation} Z-score par canal : $\tilde{x} = (x-\mu)/\sigma$, qui
        élimine les différences d'échelle entre capteurs.
  \item \textbf{Fenêtrage glissant} : découpage en fenêtres de $W=1\,024$ points avec un
        pas $H=512$ (recouvrement de 50\,\%), maximisant le nombre d'observations tout en
        préservant la résolution spectrale.
\end{enumerate}

\section{Indicateurs temporels}
Pour une fenêtre $x = (x_1,\dots,x_N)$, les descripteurs temporels calculés sont définis
par les équations~\eqref{eq:rms}--\eqref{eq:crest}. Ils caractérisent l'énergie globale
(RMS), le caractère impulsif (kurtosis, facteur de crête) et l'asymétrie (skewness) du
signal.

\begin{align}
\text{RMS} &= \sqrt{\tfrac{1}{N}\textstyle\sum_{i=1}^{N} x_i^{2}} \label{eq:rms}\\
\text{Kurtosis} &= \frac{\tfrac{1}{N}\sum_{i=1}^{N}(x_i-\mu)^4}{\sigma^{4}} - 3 \label{eq:kurt}\\
\text{Facteur de crête} &= \frac{\max_i |x_i|}{\text{RMS}} \label{eq:crest}
\end{align}

L'intérêt physique est direct : un défaut de roulement naissant produit des chocs
périodiques qui élèvent le kurtosis (au-delà de 3, valeur d'un bruit gaussien) et le facteur
de crête, \emph{avant même} que l'énergie globale (RMS) n'augmente significativement.

\section{Indicateurs fréquentiels}
La transformée de Fourier rapide (FFT) fournit le spectre $X_k$. On en extrait la fréquence
dominante, l'énergie répartie par bandes (0--1, 1--3, 3--5, 5--6{,}4\,kHz) et
l'\textbf{entropie spectrale}~\eqref{eq:ent}, qui mesure l'étalement de l'énergie : un
spectre concentré (raie nette) a une faible entropie, un spectre diffus (jeu mécanique,
bruit large bande) une entropie élevée.

\begin{equation}
H_{\text{spec}} = -\sum_k p_k \log p_k, \qquad p_k = \frac{|X_k|^2}{\sum_j |X_j|^2}
\label{eq:ent}
\end{equation}

Selon le dataset, le vecteur de features atteint 15 (CWRU, un canal), 45 (VBL, 3 canaux),
126 (CMAPSS, 21 capteurs) ou 132 features (Mechanical Faults, 4 canaux avec descripteurs
harmoniques et d'enveloppe spécialisés).

\section{Fréquences caractéristiques de roulement}
Le diagnostic des roulements repose sur le calcul de fréquences théoriques à partir de la
géométrie (nombre de billes $n$, diamètre de bille $d$, diamètre primitif $D$, angle de
contact $\phi$) et de la fréquence de rotation $f_r$ :

\begin{align}
\text{BPFO} &= \tfrac{n}{2} f_r \left(1 - \tfrac{d}{D}\cos\phi\right) \\
\text{BPFI} &= \tfrac{n}{2} f_r \left(1 + \tfrac{d}{D}\cos\phi\right) \\
\text{BSF}  &= \tfrac{D}{2d} f_r \left(1 - \left(\tfrac{d}{D}\cos\phi\right)^2\right) \\
\text{FTF}  &= \tfrac{f_r}{2} \left(1 - \tfrac{d}{D}\cos\phi\right)
\end{align}

La présence d'un pic à BPFO signe un défaut de bague externe, à BPFI un défaut de bague
interne, à BSF un défaut de bille, à FTF un défaut de cage. Le module de diagnostic recherche
ces pics dans le spectre avec une tolérance de $\pm 8\,\%$ et calcule un score de proéminence
par rapport au plancher de bruit.

\section{Analyse d'enveloppe (démodulation de Hilbert)}
Les chocs de roulement excitent des résonances haute fréquence modulées par la fréquence du
défaut. Pour les révéler, le signal est filtré passe-bande (par défaut 500--5\,000\,Hz),
puis on calcule son enveloppe via la transformée de Hilbert $\mathcal{H}$ :
\begin{equation}
e(t) = \big|\, x_{\text{bp}}(t) + j\,\mathcal{H}\{x_{\text{bp}}\}(t)\,\big|
\end{equation}
Le spectre de l'enveloppe $e(t)$ fait alors apparaître la fréquence de défaut (BPFO/BPFI/
BSF), confirmant le diagnostic issu du spectre brut.

\section{Sévérité ISO 10816 : de l'accélération à la vitesse}
Conformément à la norme, la sévérité est évaluée par la \textbf{vitesse vibratoire efficace
(mm/s)} dans la bande 10--1\,000\,Hz. Le signal d'accélération $a(t)$ est intégré en
vitesse, après filtrage passe-bande pour éliminer la dérive d'intégration :
\begin{equation}
v(t) = \int a(t)\,dt, \qquad
V_{\text{RMS}} = \sqrt{\tfrac{1}{T}\textstyle\int_0^T v(t)^2\,dt}\ \ [\text{mm/s}]
\end{equation}
La valeur $V_{\text{RMS}}$ est ensuite classée en zone A/B/C/D selon la classe de machine
(I--IV). À titre de validation, un signal d'essai de 1\,g à 50\,Hz produit
$V_{\text{RMS}} \approx 22{,}3$\,mm/s, soit la zone D (dangereux), en accord avec le calcul
analytique $V = g/(2\pi f \sqrt{2}) \times 10^3$.

\section{Indice de santé (Health Index)}
Les sorties des différents modules sont fusionnées en un \textbf{indice de santé} $HI \in
[0,100]$, combinant trois composantes pondérées : le score d'anomalie (auto-encodeur), la
RUL normalisée et la dérive vibratoire (RMS vs baseline) :
\begin{equation}
HI = 0{,}4\cdot HI_{\text{anomalie}} + 0{,}4\cdot HI_{\text{RUL}} + 0{,}2\cdot HI_{\text{vib}}
\label{eq:hi}
\end{equation}
avec $HI_{\text{anomalie}} = \max(0,\,100 - s_{\text{anomalie}})$. En l'absence d'une
composante (par exemple la RUL pour un dataset de classification), son poids est redistribué
sur les autres. Quatre seuils sémantiques structurent l'action : SAIN ($HI>70$),
SURVEILLANCE ($40$--$70$), ALERTE ($20$--$40$) et CRITIQUE ($<20$). Un lissage exponentiel
(85\,\% ancienne valeur, 15\,\% nouvelle) stabilise l'indice contre le bruit cycle à cycle.
